\begin{proposition}[Analytic Appell--Humbert boundary cocycle]
\hypertarget{tag-03IH}{}\label{wbII:prop:analytic-ah-boundary}
\StatementRecord{wbII:prop:analytic-ah-boundary}{proved}{03IH}{\SeriesRef{wbII:eq:W:5.17-A8}, \SeriesRef{wbII:eq:W:5.17av}, \SeriesRef{wbII:eq:W:5.17ax}, \SeriesRef{wbII:eq:W:5.17-A12}, \SeriesRef{wbII:eq:W:5.17-A11}, \SeriesRef{wbII:eq:W:5.17-A13}, \SeriesRef{wbII:eq:W:5.17-A14a}, \SeriesRef{wbII:eq:W:5.17-A14c}, \SeriesRef{wbII:eq:W:5.17-A14}, \SeriesRef{wbII:eq:W:5.17-A14d}, \SeriesRef{wbII:eq:W:5.17-A14e}, \SeriesRef{wbII:eq:W:5.17-A15}, \SeriesRef{wbII:eq:W:5.17-A17}, \SeriesRef{wbII:eq:W:5.17-A19a}, \SeriesRef{wbII:eq:W:5.17-A16}, \SeriesRef{wbII:eq:W:5.17-A9}, \SeriesRef{wbII:eq:W:5.17-A20}; \cite[equations~(3.3)--(3.10), pp.~3--7]{SchnellLecture3}; \cite[\S2, fundamental theorem~B, p.~80]{Cartan1957}}

Let \(P\subset B^\circ\) be the compact pair-of-pants core,
orient its boundary meridians clockwise, and retain the zero-section lifts
of \CompendiumRef{I}{inv:B3.1}.  Thus
\[
 \rho_1\rho_2\rho_0=1,
 \qquad
 \widehat\rho_1\widehat\rho_2\widetilde a_0=1,
\tag{W.93}\hypertarget{tag-03EF}{}\label{wbII:eq:W:5.17-A1}
\]
and the second relation has no terminal vertical translation.  For
\(L\in\operatorname{Pic}(X)\), let
\[
 E=c_1(L|_F)\in\operatorname{Alt}^2(\Lambda,\mathbb Z)
   =\bigwedge^2V
\]
be its fibrewise Appell--Humbert form.  Put
\(V_{\mathbb R}=\operatorname{Hom}(\Lambda,\mathbb R)\), and use
\[
 A_g=(T_g^{-1})^{\mathsf t},
 \qquad
 \langle T_g\alpha,\lambda\rangle
 =\langle\alpha,A_g^{-1}\lambda\rangle.
\tag{W.94}\hypertarget{tag-03EY}{}\label{wbII:eq:W:5.17-A2}
\]
Here $A_g$ is the linear part of the left deck action and $T_g=A_g^{-\mathsf t}$
is its coefficient action.  A clockwise loop ends at the deck translate $g$;
its endpoint return $g^{-1}$ acts on homology by $A_g^{-1}$ and on
cohomology by $A_g^{\mathsf t}=T_g^{-1}$.  The cocycle below uses $T_g$.
There is a crossed homomorphism, unique up to a principal crossed
homomorphism,
\[
 d:\pi_1(P)\longrightarrow V_{\mathbb R},
 \qquad
 d(gh)=d(g)+T_gd(h),
 \qquad
 d(g)\sim d(g)+(T_g-I)a.
\tag{W.95}\hypertarget{tag-03F9}{}\label{wbII:eq:W:5.17-A3}
\]
If \(d_i=d(\rho_i)\), then
\[
 d_1+T_1d_2+T_1T_2d_0=0,
 \qquad
 \bar d_1+\bar d_2+\bar d_0=0
 \quad\text{in }V_G\otimes\mathbb R.
\tag{W.96}\hypertarget{tag-03FM}{}\label{wbII:eq:W:5.17-A4}
\]
For the finite filling of order \(m_j\), \(j=1,2\),
\[
 N_{T_j}d_j=\iota_{v_j}E,
 \qquad
 m_j\bar d_j=[\iota_{v_j}E],
 \qquad
 N_{T_j}=I+T_j+\cdots+T_j^{m_j-1},
\tag{W.97}\hypertarget{tag-03FN}{}\label{wbII:eq:W:5.17-A5}
\]
where \((\iota_vE)(\lambda)=E(v,\lambda)\).  At the cusp,
\[
 \langle d_0,\widehat\delta\rangle=0,
 \qquad
 \bar d_0=0\quad\text{in }V_G\otimes\mathbb R.
\tag{W.98}\hypertarget{tag-03FO}{}\label{wbII:eq:W:5.17-A6}
\]
\end{proposition}



\begin{proof}

\begin{enumerate}[label=\arabic*.,leftmargin=*]
\item Choose an open Riemann-surface neighbourhood \(P^{\rm an}\Subset B^\circ\)
which deformation retracts onto \(P\), and pull the period family to its
universal cover.  Before quotienting by the fibre lattice, the resulting
space is
\[
 \widetilde{\mathcal J}
 =\widetilde P^{\rm an}\times\mathbb C^2.
\]
It is Stein and contractible, so the pullback of \(L\) is holomorphically
trivial.  In a nowhere-zero frame, write the factor belonging to a deck
transformation \(\alpha\) as
\[
 a_\alpha\in\mathcal O^*(\widetilde{\mathcal J}),
 \qquad
 a_{\alpha\beta}(x)=a_\alpha(\beta x)a_\beta(x).
\tag{W.99}\hypertarget{tag-03FP}{}\label{wbII:eq:W:5.17-A7}
\]
The space is simply connected, hence \(a_\alpha=e^{2\pi i h_\alpha}\) for
a holomorphic logarithm.  Therefore
\[
 c(\alpha,\beta)
 =h_\alpha(\beta x)+h_\beta(x)-h_{\alpha\beta}(x)\in\mathbb Z
\tag{W.100}\hypertarget{tag-03FQ}{}\label{wbII:eq:W:5.17-A8}
\]
is independent of \(x\), and associativity makes \(c\) an integral group
two-cocycle.

\item Let \(\tau_\lambda\) denote fibre translation by
\(\Pi(z)\lambda\), \(\lambda\in\Lambda\), and let \(\widehat g\) be the
zero-section lift of a base deck transformation.  Period equivariance gives
\[
 \widehat g\,\tau_\lambda\,\widehat g^{-1}
 =\tau_{A_g\lambda}.
\tag{W.101}\hypertarget{tag-03FR}{}\label{wbII:eq:W:5.17-A9}
\]
On fibre translations, the commutator of \eqref{wbII:eq:W:5.17-A8} is
\[
 E(\lambda,\mu)
 =c(\tau_\lambda,\tau_\mu)-c(\tau_\mu,\tau_\lambda).
\tag{W.102}\hypertarget{tag-03EG}{}\label{wbII:eq:W:5.17-A10}
\]
The four oriented sides of the \((\lambda,\mu)\)-parallelogram identify
this integer with the exponential connecting class of \(L|_F\); this is
the Appell--Humbert convention used in
\eqref{wbII:eq:W:5.17av}--\eqref{wbII:eq:W:5.17ax}, and agrees with the standard factor-of-automorphy
calculation
\cite[equations~(3.3)--(3.10), pp.~3--7]{SchnellLecture3}.
Conjugating the commutator by \(\widehat g\) yields
\[
 E(A_g\lambda,A_g\mu)=E(\lambda,\mu).
\tag{W.103}\hypertarget{tag-03EH}{}\label{wbII:eq:W:5.17-A11}
\]

\item Push the logarithmic central extension defined by \(c\) out along
\(\mathbb Z\hookrightarrow\mathbb R\), and denote its central element of
weight \(r\) by \(\mathbf c(r)\).  Choose lifts \(U_\lambda\) of the fibre
translations satisfying the real Appell--Humbert normalization
\[
 U_\lambda U_\mu
 =\mathbf c\!\left(\frac12E(\lambda,\mu)\right)U_{\lambda+\mu}.
\tag{W.104}\hypertarget{tag-03EI}{}\label{wbII:eq:W:5.17-A12}
\]
To see that this normalization exists, start with arbitrary lifts.  The
difference between their real cocycle and \(E/2\) has zero commutator, so
it defines an abelian extension of the free abelian group \(\Lambda\) by
\(\mathbb R\).  Choose lifts of an integral basis and extend
multiplicatively; this splits that abelian extension and gives \eqref{wbII:eq:W:5.17-A12}.
If \(U'_\lambda\) is a second normalized system, then
\(U'_\lambda U_\lambda^{-1}\) is central, say
\(\mathbf c(\langle a,\lambda\rangle)\).  Comparing the two copies of
\eqref{wbII:eq:W:5.17-A12} proves additivity in \(\lambda\), so
\(a\in V_{\mathbb R}\).  Thus every two normalized systems differ by exactly
the character change used below.

\item For \(g\in\pi_1(P)\), choose a lift \(G_g\) of \(\widehat g\) and define
\(d(g)\in V_{\mathbb R}\) by
\[
 G_gU_{A_g^{-1}\lambda}G_g^{-1}
 =\mathbf c\bigl(\langle d(g),\lambda\rangle\bigr)U_\lambda.
\tag{W.105}\hypertarget{tag-03EJ}{}\label{wbII:eq:W:5.17-A13}
\]
The exponent is additive in \(\lambda\): conjugate \eqref{wbII:eq:W:5.17-A12} and use
\eqref{wbII:eq:W:5.17-A11}.  The lifts of \(gh\) and \(G_gG_h\) differ centrally, so
\[
\begin{aligned}
 \langle d(gh),\lambda\rangle
 &=\langle d(g),\lambda\rangle
   +\langle d(h),A_g^{-1}\lambda\rangle\\
 &=\langle d(g)+T_gd(h),\lambda\rangle.
\end{aligned}
\]
This proves the crossed-homomorphism law.  Replacing
\(U_\lambda\) by
\(\mathbf c(\langle a,\lambda\rangle)U_\lambda\) changes \eqref{wbII:eq:W:5.17-A13} by
\((T_g-I)a\).  Replacing \(G_g\) by a central multiple does not change
\eqref{wbII:eq:W:5.17-A13}, so this proves the stated independence and uniqueness.  Applying
the crossed law to
\eqref{wbII:eq:W:5.17-A1} proves the first identity in \eqref{wbII:eq:W:5.17-A4}; passage to coinvariants
proves the second.

\item For the order-\(m_j\) filling, put
\[
 \widetilde J_j=\Delta_j\times\mathbb C^2,\qquad
 r_j:\widetilde J_j\longrightarrow J'_j
\]
for the lattice quotient, whose deck translations are \(\tau_\lambda\).
On the vector cover define
\[
 Q_{j,\log}:=\widetilde g_j^{\,v_j/m_j},\qquad
 Q_{j,\mathrm{lin}}:=\widetilde g_j^{\,0}.
\]
They satisfy
\[
\begin{aligned}
 Q_{j,\log}\tau_\lambda Q_{j,\log}^{-1}
  &=\tau_{A_j\lambda},&
 (Q_{j,\log})^{m_j}&=\tau_{v_j},\\
 Q_{j,\mathrm{lin}}\tau_\lambda Q_{j,\mathrm{lin}}^{-1}
  &=\tau_{A_j\lambda},&
 (Q_{j,\mathrm{lin}})^{m_j}&=1.
\end{aligned}
\tag{W.106}\hypertarget{tag-03EK}{}\label{wbII:eq:W:5.17-A14}
\]
The positive sign is fixed by
\(s_j\circ g_j=e^{-2\pi i/m_j}s_j\) and \CompendiumRef{I}{pub:I:eq:PD7}: after exactly
\(m_j\) clockwise iterations, the logarithmic vector-cover lift translates
by \(+v_j\).

\item Let \(q_{j,\log},q_{j,\mathrm{lin}}\) be the descended automorphisms of \(J'_j\),
and put \(J_j^{\prime *}=J'_j|_{\Delta_j^*}\).  Both descended automorphisms have
order dividing \(m_j\), because \(\tau_{v_j}\) is trivial after the lattice
quotient; each covers \(g_j\), whose order is exactly \(m_j\), so each
descended order is exactly \(m_j\).
The branch-independent torus translation
\[
 h_j:J_j^{\prime *}\longrightarrow J_j^{\prime *}
\]
constructed in \CompendiumRef{I}{pub:I:eq:PD8} satisfies
\[
 h_jq_{j,\log}h_j^{-1}=q_{j,\mathrm{lin}}.
\]
Define
\[
\begin{aligned}
 p_{j,\log}:J_j^{\prime *}&\longrightarrow
  N_j^*:=J_j^{\prime *}/\langle q_{j,\log}\rangle,\\
 p_{j,\mathrm{lin}}:J_j^{\prime *}&\longrightarrow
 J_j^*:=J_j^{\prime *}/\langle q_{j,\mathrm{lin}}\rangle .
\end{aligned}
\]
The conjugacy induces a biholomorphism
\[
 \bar h_j:N_j^*\longrightarrow J_j^*,\qquad
 p_{j,\mathrm{lin}}\circ h_j=\bar h_j\circ p_{j,\log}.
\]
The gluing of \(X\) identifies
\(L|_{N_j^*}\simeq\bar h_j^*(L|_{J_j^*})\); hence the
line-bundle square gives
\[
 p_{j,\log}^*(L|_{N_j^*})
 \simeq h_j^*p_{j,\mathrm{lin}}^*(L|_{J_j^*}).
\tag{W.107}\hypertarget{tag-03EL}{}\label{wbII:eq:W:5.17-A14a}
\]

\item On the filled vector cover, let
\[
 P_{j,\log}:J'_j\longrightarrow N_j,\qquad
 \pi_j^{\mathrm{vec}}:=P_{j,\log}\circ r_j:
 \widetilde J_j\longrightarrow N_j .
\]
Here \(P_{j,\log}\) is the full quotient and
\(P_{j,\log}|_{J_j^{\prime *}}=p_{j,\log}\).
The deck group is
\[
 \Gamma_j^{\log}
 =\langle\tau_\lambda,Q_{j,\log}:\lambda\in\Lambda\rangle.
\]
The pullback \((\pi_j^{\mathrm{vec}})^*(L|_{N_j})\) is trivial on the
contractible Stein manifold \(\widetilde J_j\); its factors define an
integral central extension, which we push out along
\(\mathbb Z\hookrightarrow\mathbb R\).  Over \(\Delta_j^*\), the
line-bundle square \eqref{wbII:eq:W:5.17-A14a} identifies each fibre with the global factor
extension by a translation, so its commutator form is \(E\).  That form is
integer-valued and locally constant on the connected disc \(\Delta_j\);
therefore it is \(E\) on the filled cover as well.  Choose vertical lifts
\(U_\lambda^{(j)}\) normalized by
\[
 U_\lambda^{(j)}U_\mu^{(j)}
 =\mathbf c\!\left(\frac12E(\lambda,\mu)\right)
  U_{\lambda+\mu}^{(j)}.
\]
For \(\alpha\in\Gamma_j^{\log}\), define \(A_\alpha,T_\alpha\) by
\[
 \alpha\tau_\lambda\alpha^{-1}=\tau_{A_\alpha\lambda},\qquad
 T_\alpha=(A_\alpha^{-1})^{\mathsf t}.
\]
For a lift \(\widetilde G_\alpha\), define
\(\widetilde d_j(\alpha)\in V_{\mathbb R}\) through
\[
 \widetilde G_\alpha U_{A_\alpha^{-1}\lambda}^{(j)}
 \widetilde G_\alpha^{-1}
 =\mathbf c\!\left(\langle\widetilde d_j(\alpha),\lambda\rangle\right)
  U_\lambda^{(j)}.
\tag{W.108}\hypertarget{tag-03EM}{}\label{wbII:eq:W:5.17-A14c}
\]
Indeed, before naming \(\widetilde d_j(\alpha)\), the quotient of the two
sides of \eqref{wbII:eq:W:5.17-A14c} is central because both cover the same translation.
Conjugating the normalized product and using
\[
 E(A_\alpha\lambda,A_\alpha\mu)=E(\lambda,\mu)
\]
shows that its exponent \(r_\alpha(\lambda)\) satisfies
\[
 r_\alpha(\lambda+\mu)=r_\alpha(\lambda)+r_\alpha(\mu).
\]
Thus it is the pairing with a unique element of \(V_{\mathbb R}\), as
asserted.  Since \(A_{\alpha\beta}=A_\alpha A_\beta\), choose the lift of
\(\alpha\beta\) to be \(\widetilde G_\alpha\widetilde G_\beta\); any other
choice differs centrally.  Then
\[
\begin{aligned}
 &\widetilde G_\alpha\widetilde G_\beta
 U_{A_\beta^{-1}A_\alpha^{-1}\lambda}^{(j)}
 \widetilde G_\beta^{-1}\widetilde G_\alpha^{-1}\\
 &\quad=
 \mathbf c\!\left(
  \langle\widetilde d_j(\beta),A_\alpha^{-1}\lambda\rangle
  +\langle\widetilde d_j(\alpha),\lambda\rangle\right)
 U_\lambda^{(j)}.
\end{aligned}
\]
The defining dual action
\(\langle T_\alpha v,\lambda\rangle
=\langle v,A_\alpha^{-1}\lambda\rangle\) therefore proves the crossed law
\[
 \widetilde d_j(\alpha\beta)
 =\widetilde d_j(\alpha)+T_\alpha\widetilde d_j(\beta).
\tag{W.109}\hypertarget{tag-03EN}{}\label{wbII:eq:W:5.17-A14d}
\]

\item Put
\[
 \widetilde J_j^*:=\Delta_j^*\times\mathbb C^2,
 \qquad
 r_j^*:=r_j|_{\widetilde J_j^*}:
 \widetilde J_j^*\longrightarrow J_j^{\prime *}
\]
and define
\[
\begin{aligned}
 \Phi_{j,\log}
  &:=p_{j,\log}\circ r_j^*,\\
 \Phi_{j,\mathrm{glob}}
  &:=p_{j,\mathrm{lin}}\circ h_j\circ r_j^*
   =\bar h_j\circ\Phi_{j,\log}.
\end{aligned}
\]
Both maps are invariant under \(\Gamma_j^{\log}\).  This is immediate for
every \(\tau_\lambda\), while
\[
\begin{aligned}
 \Phi_{j,\log}Q_{j,\log}
 &=p_{j,\log}q_{j,\log}r_j^*
  =\Phi_{j,\log},\\
 \Phi_{j,\mathrm{glob}}Q_{j,\log}
 &=p_{j,\mathrm{lin}}h_jq_{j,\log}r_j^*
  =p_{j,\mathrm{lin}}q_{j,\mathrm{lin}}h_jr_j^*
  =\Phi_{j,\mathrm{glob}}.
\end{aligned}
\]
Thus \(\widetilde J_j^*\) is the same \(\Gamma_j^{\log}\)-cover of the two
biholomorphic quotients.  Pulling the gluing isomorphism to this cover gives
\[
 \Theta_j:
 \Phi_{j,\log}^*(L|_{N_j^*})
 \xrightarrow{\ \sim\ }
 \Phi_{j,\mathrm{glob}}^*(L|_{J_j^*}).
\]
For the comparison of logarithms, pass one further time to the universal
cover
\[
 \widehat J_j^*=\widetilde{\Delta_j^*}\times\mathbb C^2
 \longrightarrow \widetilde J_j^*.
\]
The deck group over the punctured quotient is
\[
 \widehat\Gamma_j=
 \langle\tau_\lambda,q_j\mid
 q_j\tau_\lambda q_j^{-1}=\tau_{A_j\lambda}\rangle.
\]
The map from this universal cover to the filled vector cover induces the
explicit homomorphism
\[
 \sigma_j:\widehat\Gamma_j\longrightarrow\Gamma_j^{\log},
 \qquad \tau_\lambda\longmapsto\tau_\lambda,
 \quad q_j\longmapsto Q_{j,\log}.
\]
Its kernel is generated by
\(q_j^{m_j}\tau_{-v_j}\): this element is one deck turn of
\(\widetilde{\Delta_j^*}\to\Delta_j^*\), while its image on the
filled vector cover is the identity by \eqref{wbII:eq:W:5.17-A14}.
Indeed every word is \(\tau_\lambda q_j^r\); its filled image can be
identity only when \(r=m_jk\), and then only when
\(\lambda=-kv_j\).  Thus the kernel assertion follows directly.

\item Let \(\mathscr E_j\) be the real central extension of
\(\Gamma_j^{\log}\) defined by the filled factors.  Pull it back along
\(\sigma_j\), explicitly as the fibre-product group
\[
 \sigma_j^*\mathscr E_j=
 \{(e,g):e\in\mathscr E_j,\ g\in\widehat\Gamma_j,
           \ p(e)=\sigma_j(g)\}.
\]
The global factor extension restricted to the based boundary group is
an extension of the same \(\widehat\Gamma_j\), with the generator
\(q_j\) identified with \(\widehat\rho_j\) through the fixed whisker.
The endpoint calculation below proves this precise identification.
Pull \(\Theta_j\) to \(\widehat J_j^*\).  This space is simply
connected, so the holomorphic unit comparing the two frames has a
holomorphic logarithm.  Substitution into \eqref{wbII:eq:W:5.17-A8}
changes the two integral cocycles by a coboundary.  Consequently it gives
an isomorphism from \(\sigma_j^*\mathscr E_j\) to the restricted global
extension, sending every \(\mathbf c(r)\) to the same
\(\mathbf c(r)\).  In the comparisons below,
\(U_\lambda^{(j)}\) means its paired lift
\((U_\lambda^{(j)},\tau_\lambda)\), and
\(\widetilde G_{Q_{j,\log}}\) means
\((\widetilde G_{Q_{j,\log}},q_j)\), when the isomorphism
\(\Theta_{j,*}\) is applied.  Their conjugation exponents remain the
ones in the filled extension.  
\item Let \(\gamma_j^J:[0,1]\to J_j^*\) be the zero-section image of the
unwhiskered clockwise boundary meridian, and let \(\omega_j\) be the fixed
zero-section whisker from the global basepoint to \(\gamma_j^J(0)\).  Thus
the based total-space loop retained in \eqref{wbII:eq:W:5.17-A1} is
\[
 \widehat\rho_j=\omega_j*\gamma_j^J*\omega_j^{-1}.
\]
Let $s_j$ be the root point selected by the lift of the fixed zero-section
whisker, and choose a logarithm branch $\alpha=\log s_j/(2\pi i)$.  For the vector-cover connector
$H_j(s,\zeta)=(s,\zeta+\alpha(s)\Pi(s)v_j)$, set
\[
 \widetilde x_j=(s_j,-\alpha(s_j)\Pi(s_j)v_j),\qquad
 x_j=r_j^*(\widetilde x_j),\qquad x_{j,\mathrm{lin}}=h_j(x_j).
\]
Thus $H_j(\widetilde x_j)=(s_j,0)$ and
$p_{j,\mathrm{lin}}(x_{j,\mathrm{lin}})=\gamma_j^J(0)$.  Let
\(\widetilde\gamma_{j,\mathrm{lin}}\) be the
\(p_{j,\mathrm{lin}}\)-lift of \(\gamma_j^J\) beginning at
\(x_{j,\mathrm{lin}}\).  For the chosen clockwise generator,
\[
 \widetilde\gamma_{j,\mathrm{lin}}(1)
 =q_{j,\mathrm{lin}}(x_{j,\mathrm{lin}}).
\]
Set
\[
 \gamma_j^N:=\bar h_j^{-1}\circ\gamma_j^J,
 \qquad
 \widetilde\gamma_{j,\log}
 :=h_j^{-1}\circ\widetilde\gamma_{j,\mathrm{lin}}.
\]
Then the quotient square gives the point-valued identities
\[
\begin{aligned}
 p_{j,\log}\circ\widetilde\gamma_{j,\log}
  &=\gamma_j^N,\\
 \widetilde\gamma_{j,\log}(0)&=x_j,\\
 \widetilde\gamma_{j,\log}(1)&=q_{j,\log}(x_j).
\end{aligned}
\]
Let $\widetilde\gamma_{j,\mathrm{vec}}$ be the $r_j^*$-lift of
$\widetilde\gamma_{j,\log}$ beginning at the selected $\widetilde x_j$.
For the clockwise lift, $\alpha_1=\alpha_0-1/m_j$.  Writing
$\zeta_0=-\alpha_0\Pi(s_j)v_j$, period equivariance and $A_jv_j=v_j$ give
\[
 \zeta_1=-(\alpha_0-1/m_j)\Pi(g_js_j)v_j
 =R_{g_j}(s_j)\zeta_0+\Pi(g_js_j)v_j/m_j.
\]
By \CompendiumRef{I}{pub:I:eq:PD7}, this gives
\[
 \widetilde\gamma_{j,\mathrm{vec}}(1)
 =Q_{j,\log}(\widetilde x_j),
\]
for this selected vector lift.
Fibre loops remain \(\tau_\lambda\), since \(h_j\) is a fibre translation.

\item Lift $\omega_j$ along the zero section on the linear cover and use
$H_j^{-1}$ at its local endpoint; this gives the selected
$\widetilde x_j$ above.  The factor of $L$ along this fixed lift identifies the
target factor over the local boundary segment with the globally based
factor.  Let \(G_{\rho_j}^{(j)}\) be the target lift of the local boundary
deck element whose whiskered image is the chosen global lift
\(G_{\rho_j}\).  Thus it covers \(\widehat\rho_j\) after whiskering and,
by \eqref{wbII:eq:W:5.17-A13}, satisfies
\[
 G_{\rho_j}^{(j)}U_{A_j^{-1}\lambda}(G_{\rho_j}^{(j)})^{-1}
 =\mathbf c(\langle d_j,\lambda\rangle)U_\lambda,
 \qquad d_j=d(\rho_j).
\]
Consequently, for some
\(a_j\in V_{\mathbb R}\) and \(b_j\in\mathbb R\),
\[
 \Theta_{j,*}(U_\lambda^{(j)})
 =\mathbf c(\langle a_j,\lambda\rangle)U_\lambda,
 \qquad
 \Theta_{j,*}(\widetilde G_{Q_{j,\log}})
 =\mathbf c(b_j)G_{\rho_j}^{(j)}.
\]
Here \(U_\lambda\) is the normalized global system in \eqref{wbII:eq:W:5.17-A12}, and
\(\mathbf c(b_j)\) disappears under conjugation.
Therefore
\[
\begin{aligned}
 \langle\widetilde d_j(Q_{j,\log}),\lambda\rangle
 +\langle a_j,\lambda\rangle
 &=\langle d_j,\lambda\rangle
   +\langle a_j,A_j^{-1}\lambda\rangle\\
 &=\langle d_j+T_ja_j,\lambda\rangle,
\end{aligned}
\]
so
\[
 \widetilde d_j(Q_{j,\log})=d_j+(T_j-I)a_j.
\tag{W.110}\hypertarget{tag-03EO}{}\label{wbII:eq:W:5.17-A14e}
\]

\item For a translation choose
\(\widetilde G_{\tau_v}=U_v^{(j)}\).  The filled normalization above gives
\[
 U_v^{(j)}U_\lambda^{(j)}(U_v^{(j)})^{-1}
 =\mathbf c(E(v,\lambda))U_\lambda^{(j)},
\]
and hence
\[
 \widetilde d_j(\tau_v)=\iota_vE,\qquad
 (\iota_vE)(\lambda)=E(v,\lambda).
\tag{W.111}\hypertarget{tag-03EP}{}\label{wbII:eq:W:5.17-A14f}
\]
Here
\[
 A_{Q_{j,\log}}=A_j,\qquad T_{Q_{j,\log}}=T_j.
\]
Since \(N_{T_j}(T_j-I)=T_j^{m_j}-I=0\), equations
\eqref{wbII:eq:W:5.17-A14}, \eqref{wbII:eq:W:5.17-A14d}, and \eqref{wbII:eq:W:5.17-A14e} give
\[
\begin{aligned}
 N_{T_j}d_j
 &=N_{T_j}\widetilde d_j(Q_{j,\log})\\
 &=\widetilde d_j((Q_{j,\log})^{m_j})\\
 &=\widetilde d_j(\tau_{v_j})\\
 &=\iota_{v_j}E.
\end{aligned}
\tag{W.112}\hypertarget{tag-03EQ}{}\label{wbII:eq:W:5.17-A15}
\]
Changing a line-bundle frame gives an isomorphic factor extension; changing
\(\widetilde G_\alpha\) by a central factor does not change conjugation; and
every second normalized vertical system differs by a character, changing
\(\widetilde d_j(\alpha)\) by \((T_\alpha-I)a\).  Finally,
\(\ell\mapsto\ell+2\pi in\) changes a vector-cover lift of \(h_j\) by
\[
 H_j\longmapsto H'_j=\tau_{nv_j}H_j.
\]
At the factor-extension level this means
\[
 \Theta'_{j,*}
 =\operatorname{Ad}(U_{nv_j})\circ\Theta_{j,*}.
\]
Consequently the comparison character changes by
\(n\iota_{v_j}E\), because
\[
 U_{nv_j}
 \bigl(\mathbf c(\langle a_j,\lambda\rangle)U_\lambda\bigr)
 U_{nv_j}^{-1}
 =
 \mathbf c(\langle a_j+n\iota_{v_j}E,\lambda\rangle)U_\lambda.
\]
Since \(A_jv_j=v_j\) and \(E\) is \(A_j\)-invariant,
\[
 \langle T_j(\iota_{v_j}E),\lambda\rangle
 =E(v_j,A_j^{-1}\lambda)
 =E(A_jv_j,\lambda)
 =E(v_j,\lambda),
\]
so this branch change leaves both \eqref{wbII:eq:W:5.17-A14e} and \eqref{wbII:eq:W:5.17-A15} unchanged.
The translation \(\tau_{nv_j}\) also commutes with
\(Q_{j,\mathrm{lin}}\) on the vector cover because \(A_jv_j=v_j\); on the
lattice quotient it is already the identity.  Hence it changes only the
central lift choice there.
The descended \(h_j,\bar h_j\), and the quotient square are branch
independent.  Thus \eqref{wbII:eq:W:5.17-A15} is choice-independent and proves \eqref{wbII:eq:W:5.17-A5},
including its sign.

\item Since
\(A_0\widehat\delta=\widehat\delta\), equation \eqref{wbII:eq:W:5.17-A13} says that
\[
 [G_0,U_{\widehat\delta}]
 =\mathbf c\bigl(\langle d_0,\widehat\delta\rangle\bigr).
\tag{W.113}\hypertarget{tag-03ER}{}\label{wbII:eq:W:5.17-A16}
\]
This logarithmic commutator is the clutching winding with the base loop
first.  Thus, if \(b\) and \(f\) denote the oriented base and fibre loops,
respectively,
\[
 \operatorname{wind}_{b\wedge f}(L)
  =\langle d_0,\widehat\delta\rangle,
 \qquad
 \operatorname{wind}_{\xi_{\widehat\delta}}(L)
 =-\langle d_0,\widehat\delta\rangle.
\tag{W.114}\hypertarget{tag-03ES}{}\label{wbII:eq:W:5.17-A17}
\]
The second equality uses the fibre-first convention
\(\xi_{\widehat\delta}=[S^1_{\widehat\delta}]\times[\rho_0]\), hence the
single sign change.
Use the unimodular cusp chart of \CompendiumRef{I}{mapsI:src:A.1.2}
\[
 (u,x_1,x_2)=\left(\frac{t}{x_1x_2},x_1,x_2\right),
 \qquad t=ux_1x_2.
\tag{W.115}\hypertarget{tag-03ET}{}\label{wbII:eq:W:5.17-A18}
\]
For \(0<\rho\ll1\), the boundary parametrization
\[
 (\theta,\phi)\longmapsto
 \bigl(\rho e^{i(\theta-\phi)},1,e^{i\phi}\bigr),
 \qquad t=\rho e^{i\theta},
\tag{W.116}\hypertarget{tag-03EU}{}\label{wbII:eq:W:5.17-A19}
\]
has base coordinate \(\theta\) and fibre translation
\(\widehat\delta\) in the \(\phi\)-direction.  The latter fixed-base loop
does \emph{not} itself bound the \(x_2\)-disc.  To identify the correct
meridian, put \(\alpha=\theta-\phi\), let \(a_u\) be the positively
oriented \(u\)-longitude, and let \(m_{x_2}\) be the positively oriented
\(x_2\)-meridian at fixed \(u\).  In the basis
$ (a_u,m_{x_2}) $, the counterclockwise base loop $b_+$ and the fibre loop are
\[
 b_+=a_u,
 \qquad
 f=m_{x_2}-a_u,
 \qquad
 b_+\wedge f=a_u\wedge m_{x_2}.
\tag{W.117}\hypertarget{tag-03EV}{}\label{wbII:eq:W:5.17-A19a}
\]
The determinant of this change of basis is $+1$.  The clockwise loop in
\eqref{wbII:eq:W:5.17-A17} is $b=-b_+$, so
$b\wedge f=-a_u\wedge m_{x_2}$.  The winding computed below is zero
in either orientation.  The loop $m_{x_2}$ bounds the
analytic disc
\[
 z\longmapsto(\rho e^{i\alpha_0},1,z),
 \qquad |z|\le1.
\tag{W.118}\hypertarget{tag-03EW}{}\label{wbII:eq:W:5.17-A19b}
\]

\item Choose a thin annulus \(A_u\) about \(|u|=\rho\), a disc \(D_{x_1}\)
about \(1\), and a disc \(D_{x_2}\) containing \(\{|x_2|\le1\}\), small
enough that
\[
 A_u\times D_{x_1}\times D_{x_2}
 \subset\{|ux_1x_2|<\epsilon\}.
\tag{W.119}\hypertarget{tag-03EX}{}\label{wbII:eq:W:5.17-A19c}
\]
Write \(U=A_u\times D_{x_1}\times D_{x_2}\) for this neighbourhood.
It is Stein and retracts onto \(S^1\), so \(H^2(U,\mathbb Z)=0\).
The exponential sequence is
\[
0\longrightarrow\mathbb Z\longrightarrow\mathcal O_U
 \xrightarrow{\,h\mapsto e^{2\pi i h}\,}\mathcal O_U^*
 \longrightarrow1.
\]
Its exactness follows from the existence of a holomorphic logarithm locally
for every nonvanishing holomorphic function, and from the integral kernel
of the exponential. Cartan's theorem~B gives
\(H^1(U,\mathcal O_U)=H^2(U,\mathcal O_U)=0\)
\cite[\S2, fundamental theorem~B, p.~80]{Cartan1957}.
The long cohomology sequence therefore identifies
\(\operatorname{Pic}(U)=H^1(U,\mathcal O_U^*)\) with
\(H^2(U,\mathbb Z)=0\). Pull \(L|_{N_0}\) back to
the cusp cover and choose one holomorphic frame on this entire
neighbourhood.  Compute the clutching cocycle with this extending frame on
both sides of a cut in the \(u\)-longitude.  Its transition unit on
\(D_{x_2}\) is then identically \(F=1\), so by \eqref{wbII:eq:W:5.17-A19a} its
base--fibre commutator winding is
\[
 \frac1{2\pi i}\int_{|x_2|=1}\frac{dF}{F}=0.
\tag{W.120}\hypertarget{tag-03EZ}{}\label{wbII:eq:W:5.17-A20}
\]
This computes the same alternating logarithmic cocycle as \eqref{wbII:eq:W:5.17-A16}.
Indeed, changing from the global Appell--Humbert frame to the extending
chart frame changes \(c\) by a group coboundary.  Put
\[
 g=\widehat g_0,
 \qquad
 h=\tau_{\widehat\delta}
\]
in the underlying deck group.  Equation \eqref{wbII:eq:W:5.17-A9} and
\(A_0\widehat\delta=\widehat\delta\) give \(gh=hg\).  If
\(c'=c+\delta\beta\) is the cocycle after a frame change, then
\[
 \bigl[c'(g,h)-c'(h,g)\bigr]
 -\bigl[c(g,h)-c(h,g)\bigr]
 =\beta(hg)-\beta(gh)=0.
\]
Thus the alternating cocycle is frame-invariant without assuming that its
central-extension lifts commute.  Equations \eqref{wbII:eq:W:5.17-A17}, \eqref{wbII:eq:W:5.17-A19a}, and
\eqref{wbII:eq:W:5.17-A20} therefore give
\(\langle d_0,\widehat\delta\rangle=0\).
Finally, the matrices \(A_1,A_2\) fix
\(\widehat\delta\), so the evaluation is unchanged by every principal
change in \eqref{wbII:eq:W:5.17-A3}.  Since
\(V_G=\mathbb Z[\delta]\) and
\(\delta(\widehat\delta)=1\), its vanishing is equivalent to
\(\bar d_0=0\).

\end{enumerate}
\end{proof}
